%! Author = Saeid Yazdani
%! Date = 3/17/2024


\section{\acs{DPS} Measurements}
\label{sec:dps-measurements}

To identify the behavior of \ac{DPS} modules, a series of measurements were
taken with and without inserted \acp{DUT} in the socket.
The tester was programmed to inject known voltage levels into the load board and
the resulting signals were monitored by an oscilloscope on the test points
provided on the load board as shown in
\autoref{fig:villach-supply-analysis-probe-points}.

\begin{figure}
    \centering
    \printfitgraphics[0.58\textheight]{figures/villach-supply-analysis/probe_points}
    \caption[Probe Points on the Load board]{
        The measurements were taken at probe (test) points on the
        load board near the \acs{DUT} site.
    }
    \label{fig:villach-supply-analysis-probe-points}
\end{figure}

\subsection{\SI{1.3}{\volt} Core Voltage Bare-Board Tests}
\label{subsec:onsite-measurements-no-dut-1v3}

Bare-board tests unveiled irregular waveform shapes during supply
settling in the form of spikes after \emph{TEST START} and
\emph{TEST STOP} commands.
\autoref{fig:villach-supply-analysis-1v3-test-start-glitch} shows that
the \ac{ATE} under investigation exhibits voltage glitches shortly after
the issue of \emph{TEST START} on the \SI{1.3}{\volt} rail.

According to this measurement, the \ac{DPS} voltage
is drawn to \SI{0}{\volt} (ground) and recovers after an overshoot
peaking at nearly \SI{2}{\volt}, which is the limit set in the test procedure.
Although the glitch is occurring at a timespan of approximately
\SI{50}{\micro\second}, it may be long enough to disturb the test flow,
one reason behind unexpected reset events during the test procedure.

\begin{figure}[t]
    \centering
    \printfitgraphics[0.58\textheight]{figures/villach-supply-analysis/scope/1v3_start_glitch.jpg}
    \caption[\SI{1.3}{\volt} Rail glitch at \emph{TEST START}
    without \acs{DUT}]{
        \ac{DPS} Waveform shows glitches occurring at
        \SI{1.8}{\milli\second} after \emph{TEST START} on \SI{1.3}{\volt} rail.
    }
    \label{fig:villach-supply-analysis-1v3-test-start-glitch}
\end{figure}

For the \emph{TEST STOP} sequence, a large overshoot to \SI{6}{\volt} for a
duration of \SI{20}{\milli\second} was observed as shown in
\autoref{fig:villach-supply-analysis-1v3-test-stop-overshoot}.
Further investigations revealed that there is a high-voltage current source
discharging its stored energy into the load board after the \emph{TEST STOP}
command.
The following points can be assumed regarding the observed overshoot:

\begin{itemize}[topsep=8pt,itemsep=4pt,partopsep=4pt, parsep=4pt]

    \item The larger the on-board supply capacitance, the slower the voltage
    rise.

    \item Very large capacitors would reduce the over-voltage to harmless values.

    \item There is no dependency on any switch-off delay, the overshoot waveform
    depends solely on the event \emph{TEST STOP}.

    \item It is assumed that there is an unwanted shoot-through path from the
    power supply rails to the DPS output pin for a short time.
    Note that some internal switches short the output to ground after
    \SI{20}{\milli\second}
    (\autoref{fig:villach-supply-analysis-1v3-test-stop-overshoot}).

\end{itemize}

\begin{figure}[t]
    \centering
    \printfitgraphics[0.58\textheight]{figures/villach-supply-analysis/scope/1v3_test-stop_overshoot_6v0.jpg}
    \caption[\SI{1.3}{\volt} Rail overshoot at \emph{TEST STOP}  without
    \acs{DUT}]{
        \ac{DPS} Waveform shows a comparatively high voltage overshoot when
        issuing a \emph{TEST STOP} command on \SI{1.3}{\volt} rail.
    }
    \label{fig:villach-supply-analysis-1v3-test-stop-overshoot}
\end{figure}

\subsection{\SI{3.3}{\volt} \acs{I/O} Bare-Board Tests}
\label{subsec:onsite-measurements-no-dut-3v3}

Performing the same measurements on the \SI{3.3}{\volt} did not reveal any of
the glitches and overshoots observed on the \SI{1.3}{\volt} rail.
The supply rails were clean on both \emph{TEST START} and \emph{TEST STOP}
commands.
This supply shows excellent behavior in its waveform shape and settling speed
as shown in \autoref{fig:villach-supply-analysis-3v3-test-start}.
A minor misbehavior resembles the large over-voltage after \emph{TEST STOP}
but the amplitude of the delayed distortion is very low as captured in
\autoref{fig:villach-supply-analysis-3v3-test-stop} and does not disturb an
actual functional test.

\begin{figure}[h!]
    \centering
    \printfitgraphics[0.58\textheight]{figures/villach-supply-analysis/scope/3v3_start_zoomed_in.jpg}
    \caption[\SI{3.3}{\volt} Rail zoomed-in at \emph{TEST START}
    without \acs{DUT}]{
        \ac{DPS} Waveform shows a clean \SI{3.3}{\volt} signal with a good
        ramp-up and settling time.
    }
    \label{fig:villach-supply-analysis-3v3-test-start}
\end{figure}

\begin{figure}[h!]
    \centering
    \printfitgraphics[0.58\textheight]{figures/villach-supply-analysis/scope/3v3_test_stop.jpg}
    \caption[\SI{3.3}{\volt} Rail at \emph{TEST STOP}
    without \acs{DUT}]{
        \ac{DPS} Waveform shows the \emph{TEST STOP} on the \SI{3.3}{\volt}
        exhibiting no overshoot in comparison to the \SI{1.3}{\volt} rail.
    }
    \label{fig:villach-supply-analysis-3v3-test-stop}
\end{figure}

\subsection{\SI{5}{\volt} Mixed-Signal Bare-Board Tests}
\label{subsec:onsite-measurements-no-dut-5v0}

Measurements on the \SI{5}{\volt} supply rails revealed a similar
issue to the \SI{1.3}{\volt} rail during the \emph{TEST STOP} sequence with a
large voltage swing of \SI{5.5}{\volt} appearing at the end of the
test, after approximately \SI{10}{\milli\second}, as shown
in \autoref{fig:villach-supply-analysis-5v0-test-stop}.

\begin{figure}[h!]
    \centering
    \printfitgraphics[0.58\textheight]{figures/villach-supply-analysis/scope/5v0_stop_overshoot.jpg}
    \caption[\SI{5}{\volt} Rail at \emph{TEST STOP}
    without \acs{DUT}]{
        \ac{DPS} Waveform shows the \emph{TEST STOP} ont the \SI{5}{\volt}
        rail exhibits the same voltage swing behavior as seen in the
        \SI{1.3}{\volt} rail test.
    }
    \label{fig:villach-supply-analysis-5v0-test-stop}
\end{figure}

This observation leads to the assumption that there are unwanted switching
events happening on the load board itself, as there are multiple relays
and \ac{FET} switches that may interact with the \ac{DPS} modules due to
improper cross-connections during switch-on and switch-off events occurring on
the load board.

\subsection{\SI{1.3}{\volt} Core Voltage with \acs{DUT} Tests}
\label{subsec:onsite-meas
urements-with-dut-1v3}

Even with a \ac{DUT} inserted in the test socket, the \ac{DPS} shows the exact
same glitching behaviour in \SI{1.8}{\milli\second} after \emph{TEST START} as
shown in \autoref{fig:villach-supply-analysis-1v3-test-start-with-dut}.

\begin{figure}[h!]
    \centering
    \printfitgraphics[0.58\textheight]{figures/villach-supply-analysis/scope/1v3_test_with_dut_glitched.jpg}
    \caption[\SI{1.3}{\volt} Rail at \emph{TEST START} with \acs{DUT}]{
        With a \ac{DUT} inserted, the \ac{DPS} module exhibits a similar
        behavior compared to the test without \ac{DUT}.
    }
    \label{fig:villach-supply-analysis-1v3-test-start-with-dut}
\end{figure}

\subsection{\SI{3.3}{\volt} \acs{I/O} with \acs{DUT} Tests}
\label{subsec:onsite-measurements-with-dut-3v3}

\autoref{fig:villach-supply-analysis-3v3-test-start-with-dut} shows
the near-perfect waveform shape for the {\SI{3.3}{\volt}} supply.
Similar to the test without any \ac{DUT} inserted, the small extra voltage
ramp-up after \SI{10}{\milli\second} should not disturb test results.

\begin{figure}[h!]
    \centering
    \printfitgraphics[0.58\textheight]{figures/villach-supply-analysis/scope/3v3_test_with_dut.jpg}
    \caption[\SI{1.3}{\volt} Rail at \emph{TEST START} with \acs{DUT}]{
        \ac{DPS} Waveform on the \SI{3.3}{\volt} shows no major flaws after
        issuing \emph{TEST STOP}.
        Only a minor voltage rise can be seen.
    }
    \label{fig:villach-supply-analysis-3v3-test-start-with-dut}
\end{figure}

\subsection{\SI{5}{\volt} Mixed-Signal with \acs{DUT} Tests}
\label{subsec:onsite-measurements-with-dut-5v0}

With a \ac{DUT} inserted, the \SI{5}{\volt} shows a different behaviour compared
to the same test without a \acs{DUT}.
The ramp-up and decay glitch after \SI{10}{\milli\second} of issuing the
\emph{TEST STOP} command is still occurring, but with a much lower amplitude.
The maximum glitch voltage is \SI{1.5}{\volt} compared to the \SI{5.5}{\volt}
when there is no load inserted, as shown in
\autoref{fig:villach-supply-analysis-5v0-test-start-with-dut}.

The fact that the ramp-up has a lower magnitude when there is a load inserted
reinforces the assumption that there is a charge injection from the \ac{DPS}
internal source: The more load is connected, the lower the error voltage level.

\begin{figure}[h!]
    \centering
    \printfitgraphics[0.58\textheight]{figures/villach-supply-analysis/scope/5v0_with_dut.jpg}
    \caption[\SI{5}{\volt} Rail at \emph{TEST STOP} with \acs{DUT}]{
        \ac{DPS} Waveform of the \SI{5}{\volt} rail with \ac{DUT} inserted.
        The addition of the load has decreased the \emph{TEST STOP} ramp-up from
        \SI{5.7}{\volt} to \SI{1.5}{\volt}
    }
    \label{fig:villach-supply-analysis-5v0-test-start-with-dut}
\end{figure}

\subsection{Measurements with \acs{DUT} and \SI{10}{\micro\farad}
Bypass Capacitor}
\label{subsec:onsite-measurements-added-10uc-capacitance}

The tests on \SI{1.3}{\volt}, \SI{3.3}{\volt} and \SI{5}{\volt} rails were
repeated after insertion of a \SI{10}{\micro\farad} bypass capacitor near
the \ac{DUT} site.
The introduction of extra bypass capacitance immediately changed the shape of
\emph{TEST START} glitches and \emph{TEST STOP} ramp-up and decay.

In comparison with the same waveform without a bypass capacitor,
as shown
in \autoref{fig:villach-supply-analysis-1v3-test-start-with-dut-10uCap}, the
frequency of the glitching is lower and the undershoot of the glitch
voltage level is drastically reduced and it does not reach \SI{0}{\volt}
anymore, although the overshoot of the glitch remains the same.
It should be noted that the improvement on the glitch comes with the cost of
a very slight overshoot and undershoot at the rising edge of the rail voltage
and increased rise-time.

% FINDING
We may deduct from these findings that the control loop of the
\ac{DPS} used for the \SI{1.3}{\volt}  rail exhibits
shortcomings when dealing with capacitive loads.

For the \SI{3.3}{\volt}, the extra capacitance did not show any contribution as
this voltage rail operates flawlessly in all configurations (with and without
combinations of bypass capacitor and \ac{DUT} inserted).

\begin{figure}[h!]
    \centering
    \printfitgraphics[0.58\textheight]{figures/villach-supply-analysis/scope/1v3_start_gltich_load_and_10uCap.jpg}
    \caption[\SI{1.3}{\volt} Rail at \emph{TEST START} with \acs{DUT} and
    \SI{10}{\micro\farad} capacitor]{
        \SI{1.3}{\volt} rail at \emph{TEST START} with \acs{DUT} and
        \SI{10}{\micro\farad} bypass capacitor.
        The frequency and magnitude of the
        glitch is drastically reduced by the insertion of a bypass capacitor
        with the cost of slight overshoot and undershoot before settling and
        reduced rise time.
    }
    \label{fig:villach-supply-analysis-1v3-test-start-with-dut-10uCap}
\end{figure}

The \SI{5}{\volt} rail exhibited a leading-edge overshoot
to almost \SI{6.7}{\volt} followed by a small undershoot at the
\emph{START TEST} command before settling down to \SI{5}{\volt} indicating
a control loop related issue when dealing with capacitive loads,
as shown in \autoref{fig:villach-supply-analysis-5v0-dut-10uc} and zoomed-in
for detailed view of the overshoot and settling time in
\autoref{fig:villach-supply-analysis-5v0-dut-10uc-zoom}.
The addition of the capacitor did affect the ramp-up voltage after issuing
\emph{TEST STOP} command.

\begin{figure}[h!]
    \centering
    \printfitgraphics[0.58\textheight]{figures/villach-supply-analysis/scope/5v0_test-start_overshoot_zoomout_dut-10uc.jpg}
    \caption[\SI{5.0}{\volt} Rail at \emph{TEST START} with \acs{DUT} and
    \SI{10}{\micro\farad} capacitor]{
        \SI{5}{\volt} rail with \acs{DUT} and
        \SI{10}{\micro\farad} bypass capacitor inserted.
        A spike to \SI{6.7}{\volt} is observed before settling down at
        \SI{5}{\volt}.

    }
    \label{fig:villach-supply-analysis-5v0-dut-10uc}
\end{figure}

\begin{figure}[h!]
    \centering
    \printfitgraphics[0.58\textheight]{figures/villach-supply-analysis/scope/5v0_start_dut_10uC.jpg}
    \caption[Detailed view of \SI{5.0}{\volt} rail at \emph{TEST START}]{
        Zooming in the \SI{5}{\volt} rail with \acs{DUT} and
        \SI{10}{\micro\farad} bypass capacitor inserted.
        The overshoot and undershoot before settling at the target voltage are
        visible at this time scale.
    }
    \label{fig:villach-supply-analysis-5v0-dut-10uc-zoom}
\end{figure}

The captured result of the  \SI{5}{\volt} with \ac{DUT} and bypass capacitor
verifies that the extra ramp-up voltage after issuing a \emph{TEST STOP} command
is indeed a charge injection phenomenon, indicating that the \ac{DPS} is unfit
for capacitive loads.

\subsection{Measurements with \acs{DUT} and \SI{20}{\micro\farad}
Bypass Capacitor}
\label{subsec:onsite-measurements-added-20uc-capacitance}

To confirm the weakness of \ac{DPS} units for capacitive loads, a
\SI{20}{\micro\farad} bypass capacitor was used and the measurements were
repeated.
The reaction of \SI{1.3}{\volt} at \emph{TEST START} rail remained overall the
same in comparison with \SI{10}{\micro\farad} capacitor and only the
overshoot and undershoot of the leading edge increased slightly before finally
settling down.

On the other hand, the \emph{TEST STOP} command exhibited a massive
spike of \SI{5.4}{\volt} that took \SI{20}{\milli\second} to disappear.
Since the time base of this measurement was set to
\SI{50}{\milli\second}, a second ramp-up and decay was observed as well.
This is shown in \autoref{fig:villach-supply-analysis-1v3-dut-20uc-stop}.

The \SI{3.3}{\volt} rail which did not exhibit any
glitch, overshoot, or undershoot in the previous measurement cases, showed an
overshoot at the \emph{TEST START} command with an amplitude of \SI{4.2}{\volt}
and slightly longer rise time,
as shown in \autoref{fig:villach-supply-analysis-3v3-dut-20uc-zoom}.

\begin{figure}[h!]
    \centering
    \printfitgraphics[0.58\textheight]{figures/villach-supply-analysis/scope/1v3_dut_20uf_50ms.jpg}
    \caption[\SI{1.3}{\volt} Rail at \emph{TEST STOP} with \acs{DUT}
    and \SI{20}{\micro\farad} capacitor]{
        The \SI{3.3}{\volt} rail with \acs{DUT} and
        \SI{20}{\micro\farad} bypass capacitor inserted.
        An extreme ramp-up to \SI{5.4}{\volt} and
        a following ramp-up is observed.
    }
    \label{fig:villach-supply-analysis-1v3-dut-20uc-stop}
\end{figure}

\begin{figure}[h!]
    \centering
    \printfitgraphics[0.58\textheight]{figures/villach-supply-analysis/scope/3v3_start_dut_20uf.jpg}
    \caption[\SI{3.3}{\volt} Rail at \emph{TEST START} with \acs{DUT}
    and \SI{20}{\micro\farad} capacitor]{
        Zooming in the \SI{3.3}{\volt} rail with \acs{DUT} and
        \SI{20}{\micro\farad} bypass capacitor inserted.
        An overshoot up to \SI{4.2}{\volt} and
        slightly slower rise-time is observed.
    }
    \label{fig:villach-supply-analysis-3v3-dut-20uc-zoom}
\end{figure}

